% !TeX encoding = UTF-8
\documentclass[variant=guia,cover=institutional,mode=screen]{ua-engineering-v6-article}
% !TeX encoding = UTF-8
\UASetUniversity{Universidad Austral}
\UASetFaculty{Facultad de Ingeniería}
\UASetDepartment{Departamento de Computación}
\UASetCampus{Pilar}
\UASetCareer{Ingeniería Informática}
\UASetCommission{Comisión A}
\UASetProfessor{Equipo docente}
\UASetSemester{Segundo cuatrimestre 2026}
\UASetCourse{Sistemas Distribuidos}
\UASetDate{2026}

\UASetSemester{Segundo cuatrimestre 2026}
\UASetDate{2026}

\UASetClassNumber{06}
\UASetClassTitle{Storage distribuido y almacenamiento permanente}
\UASetDocKind{Guía de ejercicios}
\title{Clase 06 --- Guía de ejercicios: persistencia, concurrencia y rendimiento bajo fallas}
\author{\UAProfessor}
\date{\UADate}

\begin{document}
\maketitle

\section{Propósito y método de resolución}

La guía sigue la misma progresión de las diapositivas y los apuntes:

\[
\text{contexto} \rightarrow \text{intuición} \rightarrow \text{regla}
\rightarrow \text{falla} \rightarrow \text{evidencia} \rightarrow \text{decisión}.
\]

No se evalúa la cantidad de nombres de productos. En cada respuesta se espera identificar, cuando corresponda:

\begin{enumerate}
\item el problema concreto;
\item la falla considerada;
\item la garantía buscada;
\item el mecanismo y su frontera;
\item el costo o trade-off;
\item la evidencia capaz de validar o refutar la decisión.
\end{enumerate}

\begin{uakeybox}
Una respuesta como ``usar WAL'', ``agregar un índice'' o ``elegir RocksDB'' es insuficiente si no explica qué comportamiento sostiene, ante qué falla, con qué costo y mediante qué evidencia.
\end{uakeybox}

\section{Ruta recomendada de uso}

\begin{itemize}
\item \textbf{Durante la clase:} ejercicios 1, 4, 5, 11, 15 y 19.
\item \textbf{Consolidación posterior:} ejercicios 2, 3, 6, 8, 12, 13, 16, 17 y 18.
\item \textbf{Desafío y Trabajo Final:} ejercicios 7, 9, 10, 14 y 20.
\end{itemize}

Los ejercicios 5 y 11 deben resolverse mediante timelines. Los ejercicios 15 y 19 deben apoyarse en evidencia: plan de ejecución o métricas. El ejercicio 20 exige una defensa arquitectónica completa.

\section{Nivel 1: contexto, vocabulario y promesa de negocio}

\begin{uaexercise}
\textbf{1. Del problema general a los mecanismos.} Para CampusShop, explique por qué un sistema de almacenamiento necesita cinco responsabilidades: registrar decisiones, controlar acceso concurrente, elegir rutas físicas, optimizar escrituras y proteger capacidad. Para cada responsabilidad, proponga una falla o degradación observable si no estuviera cubierta.

\textbf{Producto esperado:} tabla de cinco filas con responsabilidad, escenario, garantía y mecanismo candidato.
\end{uaexercise}

\begin{uaexercise}
\textbf{2. Siglas y términos sin ambigüedad.} Defina en español y en una frase técnica: WAL, MVCC, LSM-tree, ACK, LSN, SQL, DDL, OLTP, SLO, RPO y RTO. Agregue definiciones para COMMIT, log durable, flush durable, snapshot, serializabilidad, SSTable, merge ordenado, filtro de Bloom y write stall.

\textbf{Condición:} ninguna definición puede utilizar otra sigla no explicada.
\end{uaexercise}

\begin{uaexercise}
\textbf{3. Identidad lógica y objetos físicos.} La compra con ID 9F2A inserta una fila en \texttt{purchases}, descuenta una unidad en \texttt{inventory} y actualiza un índice. Distinga:

\begin{enumerate}
\item operación lógica;
\item fila y tabla;
\item página física;
\item entrada de índice;
\item registro WAL;
\item resultado idempotente.
\end{enumerate}

Explique qué información debe persistirse atómicamente para que un reintento con 9F2A no duplique la compra.
\end{uaexercise}

\begin{uaexercise}
\textbf{4. Prediagnóstico razonado.} Responda sin comenzar por el nombre de un mecanismo:

\begin{enumerate}
\item La aplicación mostró ``compra confirmada'' y se cortó la energía. ¿Qué debe ocurrir al reiniciar?
\item Un reporte comenzó antes de 900 compras nuevas. ¿Qué comportamiento sería coherente?
\item Para obtener 20 filas entre un millón, ¿qué trabajo físico sería deseable evitar?
\item Ingresan 90 MB/s y el mantenimiento procesa 60 MB/s. ¿Qué evolución espera?
\item Tres réplicas conservan el dato solo en RAM. ¿Qué falla común sigue sin tolerarse?
\end{enumerate}

Después asocie cada respuesta con el mecanismo estudiado en clase.
\end{uaexercise}

\section{Nivel 2: WAL, COMMIT y recuperación}

\begin{uaexercise}
\textbf{5. Timeline WAL y prefijo durable.} Dibuje cinco carriles: cliente, servicio, memoria del motor, WAL y páginas. Ubique los siguientes eventos:

\begin{enumerate}
\item solicitud 9F2A;
\item preparación de cambios;
\item append de registros \texttt{WAL UPDATE};
\item páginas sucias;
\item append del registro \texttt{COMMIT};
\item flush del WAL hasta \texttt{commitLSN};
\item ACK;
\item crash;
\item REDO.
\end{enumerate}

Responda con precisión: ¿el flush durable comienza cuando se agrega el primer \texttt{WAL UPDATE}? ¿Qué rango del log debe quedar durable antes del ACK?
\end{uaexercise}

\begin{uaexercise}
\textbf{6. Dos reglas, dos fallas.} Explique por separado:

\begin{enumerate}
\item la regla ``log antes que página'';
\item la regla ``COMMIT antes que ACK''.
\end{enumerate}

Construya un contraejemplo para cada regla invirtiendo su orden y muestre qué garantía se rompe.
\end{uaexercise}

\begin{uaexercise}
\textbf{7. Matriz de crashes y reconstrucción completa.} Considere cuatro puntos de caída:

\begin{enumerate}
\item antes del registro de COMMIT;
\item después del append de COMMIT, pero antes del flush durable;
\item después del flush durable y antes del ACK;
\item después del ACK y antes de escribir la página P17.
\end{enumerate}

Para cada caso indique qué puede observar el cliente, qué debe hacer recovery y si la compra debe aparecer confirmada. Después responda una objeción: si el COMMIT es durable, pero las páginas de \texttt{purchases}, \texttt{inventory} y del índice todavía contienen el estado anterior, ¿de dónde obtiene recovery la información necesaria para reconstruir la data completa? Distinga el papel de las páginas persistentes, los registros \texttt{WAL UPDATE} y el registro \texttt{COMMIT}.
\end{uaexercise}

\begin{uaexercise}
\textbf{8. REDO repetible.} Una página P17 conserva \texttt{pageLSN}=120. Recovery procesa registros con LSN 118, 120, 121 y 123. Proponga una regla para decidir cuáles reaplicar. Luego explique por qué esa regla permite reiniciar recovery después de otro crash.

\textbf{Extensión:} distinga idempotencia física de REDO de idempotencia de negocio mediante 9F2A.
\end{uaexercise}

\begin{uaexercise}
\textbf{9. Group commit cuantitativo.} Un \texttt{fsync} cuesta 2,4 ms. Llegan 600 transacciones por segundo. Compare:

\begin{enumerate}
\item un \texttt{fsync} por transacción;
\item grupos de 12 transacciones con una espera máxima de 1 ms para formar el grupo.
\end{enumerate}

Calcule sincronizaciones por segundo, costo físico amortizado por transacción y discuta la latencia adicional. Explique qué condición debe cumplirse para no debilitar la garantía de ACK.
\end{uaexercise}

\begin{uaexercise}
\textbf{10. Checkpoint, backup y réplica.} Un motor genera 180 MB/s de WAL. Compare checkpoints cada 2, 5 y 15 minutos. Calcule el WAL máximo posterior al último checkpoint en cada caso. Después explique qué problema resuelve un checkpoint y cuáles siguen requiriendo backup o replicación.

\textbf{Producto esperado:} tabla con volumen, posible impacto en RTO y costo operativo.
\end{uaexercise}

\section{Nivel 3: MVCC, visibilidad y serializabilidad}

\begin{uaexercise}
\textbf{11. Timeline MVCC.} Dibuje cuatro carriles: lector $T_R$, escritor $T_W$, versiones y lectura visible. Ubique:

\begin{enumerate}
\item $v_1=120$ confirmada;
\item inicio de $T_R$ y snapshot $S_R$;
\item primera lectura de $v_1$;
\item creación de $v_2=135$;
\item COMMIT de $T_W$;
\item segunda lectura dentro de $T_R$;
\item lectura de una transacción nueva;
\item fin de $T_R$ y momento en que $v_1$ queda reclamable.
\end{enumerate}

Justifique cada versión visible sin usar ``porque MVCC lo hace''.
\end{uaexercise}

\begin{uaexercise}
\textbf{12. Snapshot y vista coherente.} Explique por qué un snapshot no es una copia física completa. Proponga una regla simplificada de visibilidad para una versión creada por $T_c$ y reemplazada por $T_r$. Luego dé un ejemplo de resultado incoherente que surgiría si cada fila usara una regla temporal diferente.
\end{uaexercise}

\begin{uaexercise}
\textbf{13. Transacción larga y bloat.} Una transacción de reporte permanece abierta 90 minutos mientras se actualizan 30.000 filas por minuto. Explique:

\begin{enumerate}
\item por qué puede retener versiones anteriores;
\item qué señales operativas esperaría observar;
\item qué riesgos aparecen para tabla, índices, caché y recuperación;
\item qué intervenciones consideraría sin matar arbitrariamente la consulta.
\end{enumerate}
\end{uaexercise}

\begin{uaexercise}
\textbf{14. MVCC no implica serializabilidad.} Construya un write skew con dos guardias o dos cuentas. Dibuje snapshots, lecturas y escrituras. Explique por qué no existe conflicto de escritura sobre la misma fila y, sin embargo, se viola un invariante. Compare tres mitigaciones.
\end{uaexercise}

\section{Nivel 4: índices y lectura de evidencia}

\begin{uaexercise}
\textbf{15. De la visibilidad a la localización: árbol B y planes.} Analice la consulta:

\begin{verbatim}
SELECT id, total, confirmed_at
FROM purchases
WHERE student_id = 8421
  AND status = 'CONFIRMED'
ORDER BY confirmed_at DESC
LIMIT 20;
\end{verbatim}

Primero explique por qué MVCC y un índice responden preguntas distintas: MVCC decide qué versión puede observar la transacción; el índice reduce dónde buscar las versiones candidatas. Luego dibuje una ruta conceptual por un árbol B desde la raíz hasta una hoja y marque dónde se aplica la verificación de visibilidad.

Finalmente, compare dos planes: A muestra \texttt{Seq Scan}, 9.980 filas descartadas, 840 buffers y un \texttt{Sort}; B muestra \texttt{Index Scan}, 20 filas y 5 buffers, sin \texttt{Sort}. Léalo de abajo hacia arriba y explique qué evidencia justifica el cambio.

\textbf{Advertencia:} trate los tiempos como ilustrativos, no como promesa universal.
\end{uaexercise}

\begin{uaexercise}
\textbf{16. Diseñar y cuestionar un índice compuesto.} Escriba la sentencia DDL para crear un índice sobre \texttt{(student\_id, status, confirmed\_at DESC)}. Explique por qué ese orden sirve a la consulta anterior y cómo el alto \emph{fan-out} de un árbol B permite mantener poca altura aun con millones de entradas.

Analice estas variantes:

\begin{enumerate}
\item consulta solo por \texttt{status};
\item consulta por rango de fecha para todos los estudiantes;
\item consulta por estudiante sin estado;
\item tabla donde 95\% de las filas están confirmadas.
\end{enumerate}

Para cada variante indique si reutilizaría el índice, diseñaría otro o aceptaría un scan secuencial, y qué evidencia pediría a \texttt{EXPLAIN ANALYZE}.
\end{uaexercise}

\begin{uaexercise}
\textbf{17. El índice también cuesta.} Una tabla recibe 8.000 inserciones por segundo y posee seis índices secundarios. Proponga un experimento para decidir si dos índices poco usados deben eliminarse. Incluya métricas de lectura, escritura, WAL, espacio, cache y planes. Defina una estrategia de rollback.
\end{uaexercise}

\section{Nivel 5: LSM, filtro de Bloom, compactación y stalls}

\begin{uaexercise}
\textbf{18. Ruta completa de una clave en una LSM-tree.} Para la clave \texttt{sensor-17@12:05:03}, dibuje:

\begin{enumerate}
\item camino de escritura: WAL, memtable, memtable inmutable, flush y SSTable L0;
\item camino de lectura: memtables, SSTables candidatas, filtro de Bloom, índice, caché y bloque;
\item aparición de una versión nueva y de un tombstone;
\item condición para eliminar de forma segura la versión vieja y la marca de borrado.
\end{enumerate}

Explique por qué ``archivo imposible'' es una expresión incorrecta y reemplácela por una formulación precisa.
\end{uaexercise}

\begin{uaexercise}
\textbf{19. Diagnóstico cuantitativo de compactación.} Un dashboard muestra:

\begin{itemize}
\item ingreso lógico: 90 MB/s;
\item capacidad de compactación: 60 MB/s;
\item backlog actual: 18 GB;
\item umbral de slowdown: 27 GB;
\item umbral de stop: 36 GB;
\item archivos L0: 28 de un stop trigger de 36;
\item p99 de escritura creciente.
\end{itemize}

Calcule el déficit sostenido y el tiempo ideal hasta cada umbral si nada cambia. Construya al menos cuatro hipótesis causales y, para cada una, indique una métrica que la apoyaría o refutaría. Explique cuándo un stall es protección y cuándo evidencia de incapacidad.
\end{uaexercise}

\section{Nivel 6: diseño integral y defensa}

\begin{uaexercise}
\textbf{20. Dos cargas, dos arquitecturas defendibles.} Diseñe una estrategia de almacenamiento para:

\begin{enumerate}
\item \textbf{CampusShop}: compras, inventario, idempotencia y consultas por estudiante/fecha;
\item \textbf{telemetría del campus}: 80.000 eventos/s, rangos por dispositivo/tiempo y retención temporal.
\end{enumerate}

Para cada caso aplique $D=(P,F,G,S,T)$ e incluya:

\begin{itemize}
\item patrón de lectura y escritura;
\item modelo de falla;
\item definición de durabilidad;
\item mecanismo de recovery;
\item estrategia de concurrencia;
\item índice o estructura física;
\item política de mantenimiento;
\item SLO, RPO y RTO;
\item experimento que podría refutar la elección;
\item condición explícita para revisarla.
\end{itemize}
\end{uaexercise}

\section{Rúbrica compacta}

\begin{center}
\begin{tabular}{p{0.19\textwidth}p{0.20\textwidth}p{0.23\textwidth}p{0.25\textwidth}}
\toprule
Dimensión & Inicial & Logrado & Avanzado \\
\midrule
Contexto y falla & Enumera componentes & Declara operación y falla & Compara fallas y supuestos alternativos \\
Garantía & Usa palabras vagas & Formula una propiedad observable & Delimita su alcance y contraejemplo \\
Mecanismo & Nombra tecnología & Explica orden, versión o estructura & Justifica alternativa y frontera \\
Costo & No aparece & Identifica un trade-off & Cuantifica y conecta con SLO \\
Evidencia & Opinión & Métrica o timeline & Experimento capaz de refutar \\
\bottomrule
\end{tabular}
\end{center}

\end{document}
